\section{Analysis: sheaves on compact spaces and Banach lattices}\label{sec:analysis}

The decks so far were presheaf topoi. For analysis, the natural decks are topoi of \emph{sheaves} on a topological space $X$: a sheaf assigns a set $F(U)$ to each open set, with restrictions, such that compatible local sections glue uniquely (continuous functions, solutions of a differential equation, sections of a bundle). Negation is computed on open sets: the pseudocomplement of an open set $U$ is the interior of its complement. Two topological arches are classical:
\begin{itemize}
\item $\Sh(X)$ is Boolean iff every open subset of $X$ is closed;
\item $\Sh(X)$ is De Morgan iff $X$ is \emph{extremally disconnected}: the closure of every open set is open \cite{johnstone-demorgan}.
\end{itemize}
The second registry is functional analysis: by Gelfand duality, a compact Hausdorff space $X$ is the same as the commutative C$^*$-algebra $C(X)$, or the Banach lattice $C(X,\mathbb R)$.

\begin{theorem}\label{th:analysis}
Let $X$ be a compact Hausdorff space.
\begin{enumerate}
\item $\Sh(X)$ is Boolean iff $X$ is finite, iff $C(X)$ is finite-dimensional.
\item $\Sh(X)$ is De Morgan iff the lattice $C(X,\mathbb R)$ is \emph{Dedekind complete}: every set of continuous functions bounded above has a least upper bound among continuous functions.
\item If $X$ is metrisable, $\Sh(X)$ is De Morgan iff $X$ is finite. In particular the topoi of sheaves on $[0,1]$, on the circle, on the Cantor set are not De Morgan, while the topoi of sheaves on the Stone--Čech compactification $\beta\mathbb N$ and on the Gelfand spectrum of $L^\infty[0,1]$ are De Morgan and not Boolean.
\end{enumerate}
\end{theorem}

\begin{trace}
\item[deck] $\Sh(X)$, $X$ compact Hausdorff, with $X = $ the Gelfand spectrum of $C(X)$.
\item[arch] Bridge~\ref{br:bool} and its topological reading (open sets closed); the De Morgan arch \cite{johnstone-demorgan}.
\item[specialise] for compact Hausdorff $X$: every open set closed $\iff$ $X$ discrete $\iff$ $X$ finite; extremally disconnected $\iff$ $C(X,\mathbb R)$ Dedekind complete (\cite[Thm~14]{stone-bounded}, and Nakano \cite{nakano}); a compact metrisable extremally disconnected space is finite.
\item[rewrite] state everything on $C(X)$.
\item[verify] direct proof below; not in the Lean library (Mathlib has \texttt{ExtremallyDisconnected} and \texttt{Stonean}, but not $\Sh(X)$ as a topos with its negation).
\item[novelty] the internal form is in print: Johnstone shows that $\Sh(X)$ is De Morgan iff its object of Dedekind reals is internally conditionally order complete \cite{johnstone-demorgan}; the external form (2) via $C(X,\mathbb R)$, combining it with Stone and Nakano \cite{stone-bounded,nakano}, and the metrisable case (3), not found.
\end{trace}

\direct{} (1) If every open set is closed, every singleton (closed, as $X$ is Hausdorff) has an open complement which is closed, so singletons are open; $X$ is discrete and compact, hence finite. Then $C(X)\cong\mathbb C^{X}$. (2) This is the composite of the two readings: $\Sh(X)$ De Morgan $\iff$ $X$ extremally disconnected \cite{johnstone-demorgan} $\iff$ $C(X,\mathbb R)$ Dedekind complete \cite{stone-bounded}. The converse half of Stone's theorem is easy to see: for an open set $U$, any continuous upper bound of the functions $0\le g\le 1$ supported in $U$ is $\ge 1$ on $\overline U$ (Urysohn), and for $x\notin\overline U$ some such upper bound vanishes at $x$; so a least upper bound in $C(X,\mathbb R)$ must be the indicator function of $\overline U$, which is continuous only if $\overline U$ is open. (3) In a metrisable extremally disconnected space every convergent sequence of distinct points $x_n\to x$ gives disjoint open sets $U$, $V$ around $\{x_{2n}\}$ and $\{x_{2n+1}\}$ whose closures both contain $x$, contradicting the fact that disjoint open sets have disjoint closures in such a space; so a compact metrisable one has no non-trivial convergent sequences and is finite. $\beta\mathbb N$ and the spectrum of $L^\infty[0,1]$ are extremally disconnected (the latter because $L^\infty[0,1]$ is Dedekind complete), and infinite.\qed

\inwords{} The intuitionistic logic of \q{continuous functions on $X$} satisfies $\neg p\vee\neg\neg p$ exactly when suprema of continuous functions stay continuous. On $[0,1]$ it fails: the functions $\min(1, n\cdot\max(0, x-\tfrac12))$ are bounded by $1$ and their supremum, the indicator of $(\tfrac12,1]$, is not continuous; on the spectrum of $L^\infty$, where every bounded family has an essential supremum, it holds.

\subsection{One space, three classical theorems}

The Morita arch gives a classical equivalence of the kind discussed in \S\ref{sec:equivalences}. A sober space is determined by its topos of sheaves: $\Sh(X)\simeq\Sh(Y)$ iff $X\cong Y$, the points of $\Sh(X)$ being the points of $X$. For compact Hausdorff spaces (which are sober), read the same invariant, the space of points, on three presentations:
\begin{itemize}
\item on $\Sh(X)$: two compact Hausdorff spaces with equivalent topoi of sheaves are homeomorphic;
\item on the ring $C(X)$: the \emph{Gelfand--Kolmogorov theorem}, $C(X)\cong C(Y)$ as rings iff $X\cong Y$ (points are maximal ideals);
\item on the Banach space $C(X)$: the \emph{Banach--Stone theorem}, $C(X)$ and $C(Y)$ linearly isometric iff $X\cong Y$ (points are extreme points of the dual ball, up to sign).
\end{itemize}
The three statements are classical and proved in three different courses (Gelfand--Kolmogorov 1939; Banach 1932 and Stone 1937; the reconstruction of sober spaces from their topoi); the bridge says they are one statement, \q{the space of points is an invariant}, read on three presentations of the same object. We did not find them presented together.
